\problema{Word Search}{79}{src/01-grafos/79-word-search.cpp}{M}

Nos dan una cuadrícula de letras (\texttt{board}) y una palabra
(\texttt{word}). Queremos saber si la palabra se puede formar recorriendo
celdas adyacentes en horizontal o vertical (nunca en diagonal), sin usar la
misma celda dos veces dentro de un mismo camino.

\paragraph{Intuición.} Piensa que vas letra por letra de la palabra, parado
en una celda, y buscas a cuál vecina moverte para encontrar la siguiente
letra. Si te atoras, retrocedes (\emph{backtracking}) y pruebas otro camino.
Para no pisar dos veces la misma celda en el mismo intento, la marcamos
temporalmente (por ejemplo con \texttt{'\#'}) mientras la estamos usando, y
la restauramos a su letra original en cuanto retrocedemos, de modo que otros
caminos sí puedan usarla.

\begin{center}
\ttfamily
\begin{tabular}{cccc}
A & B & C & E\\
S & F & C & S\\
A & D & E & E\\
\end{tabular}
\end{center}

\noindent En esta cuadrícula, \texttt{"ABCCED"} existe (bajando por la
columna central de la \texttt{C} y luego a la \texttt{E} y la \texttt{D}), y
\texttt{"SEE"} también existe. En cambio \texttt{"ABCB"} no existe: tras usar
la primera \texttt{B}, la única forma de volver a una \texttt{B} obligaría a
reutilizar una celda ya visitada en ese mismo camino.

\paragraph{Solución.} Probamos cada celda de la cuadrícula como posible
punto de partida (\texttt{word[0]}) y hacemos DFS con backtracking: si la
celda actual coincide con la letra esperada, la marcamos como visitada y
probamos las 4 vecinas buscando la siguiente letra; si ningún vecino
completa el camino, restauramos la celda antes de volver.

\lstinputlisting{src/01-grafos/79-word-search.cpp}

\complejidad{$O(m\cdot n\cdot 4^{L})$, con $L$ la longitud de \texttt{word}}{$O(L)$ por la recursión}

\paragraph{Notas de entrevista (Amazon).} El punto clave es el backtracking:
marcar la celda al entrar y \textbf{siempre} restaurarla al salir, incluso
si el camino tuvo éxito antes de regresar la llamada. Si no puedes mutar la
cuadrícula de entrada, usa una matriz \texttt{visited} paralela en vez de
\texttt{'\#'}. Casos borde: palabra de una sola letra (basta con que exista
en la cuadrícula), cuadrícula de una sola celda, y palabras que reutilizarían
una celda (deben fallar).
